\chapter{Cơ sở Lý thuyết}

\section{Mạng Nơ-ron Tích chập Đồ thị (GCN)}

\subsection{Lý thuyết Đồ thị và Biểu diễn Mạng lưới Giao thông}

Mạng lưới đường phố được mô hình hóa tự nhiên bằng \textbf{đồ thị có hướng} $\mathcal{G} = (\mathcal{V}, \mathcal{E}, \mathbf{W})$, trong đó:

\begin{itemize}[itemsep=3pt]
  \item $\mathcal{V} = \{v_1, v_2, \ldots, v_N\}$: Tập hợp $N$ \textbf{đỉnh} (nodes) --- mỗi đỉnh đại diện cho một nút giao hoặc đoạn đường có gắn camera.
  \item $\mathcal{E} \subseteq \mathcal{V} \times \mathcal{V}$: Tập hợp các \textbf{cạnh} (edges) --- mỗi cạnh $(v_i, v_j)$ biểu diễn đoạn đường nối trực tiếp từ nút $v_i$ đến nút $v_j$.
  \item $\mathbf{W} \in \mathbb{R}^{N \times N}$: \textbf{Ma trận trọng số} --- $W_{ij}$ phản ánh mức độ ảnh hưởng giao thông từ nút $i$ đến nút $j$, thường được tính bằng hàm giảm theo khoảng cách hoặc thời gian di chuyển.
\end{itemize}

Trọng số thường được thiết kế theo công thức Gaussian:
\begin{equation}
  W_{ij} = \exp\!\left(-\frac{d_{ij}^2}{\sigma^2}\right)
  \label{eq:edge_weight}
\end{equation}
với $d_{ij}$ là khoảng cách (km) hoặc thời gian di chuyển (phút) giữa hai nút, $\sigma$ là tham số điều chỉnh độ rộng.

\subsection{Phép Tích chập Trên Đồ thị}

Khác với ảnh 2D có lưới đều đặn, dữ liệu trên đồ thị không có cấu trúc lưới. Phép tích chập đồ thị (Graph Convolution) được định nghĩa trong miền phổ (spectral domain) thông qua phân tích riêng (eigen-decomposition) của ma trận Laplacian chuẩn hóa:

\begin{equation}
  \mathbf{L} = \mathbf{I}_N - \mathbf{D}^{-1/2}\mathbf{A}\mathbf{D}^{-1/2}
  \label{eq:laplacian}
\end{equation}

Kipf \& Welling \cite{kipf2017gcn} đề xuất xấp xỉ bậc nhất (first-order approximation) rất hiệu quả:

\begin{equation}
  \mathbf{H}^{(l+1)} = \sigma\!\left(\tilde{\mathbf{D}}^{-1/2}\tilde{\mathbf{A}}\tilde{\mathbf{D}}^{-1/2}\mathbf{H}^{(l)}\mathbf{W}^{(l)}\right)
  \label{eq:gcn_update}
\end{equation}

trong đó:
\begin{itemize}[itemsep=3pt]
  \item $\tilde{\mathbf{A}} = \mathbf{A} + \mathbf{I}_N$: Ma trận kề có thêm self-loop (để mỗi nút nhận thông tin từ chính nó).
  \item $\tilde{\mathbf{D}}_{ii} = \sum_j \tilde{A}_{ij}$: Ma trận bậc (degree matrix) của $\tilde{\mathbf{A}}$.
  \item $\mathbf{H}^{(l)} \in \mathbb{R}^{N \times F_l}$: Ma trận đặc trưng ẩn tại lớp $l$, với $F_l$ chiều đặc trưng.
  \item $\mathbf{W}^{(l)} \in \mathbb{R}^{F_l \times F_{l+1}}$: Ma trận trọng số học được của lớp $l$.
  \item $\sigma(\cdot)$: Hàm kích hoạt phi tuyến (ReLU, LeakyReLU).
\end{itemize}

\begin{tcolorbox}[tipbox]
\textbf{Ý nghĩa vật lý:} Phép tích chập đồ thị cho phép mô hình học cách \textit{ùn tắc lan truyền}: nếu nút giao A bắt đầu tắc, sau một lớp GCN, thông tin này được truyền sang các nút lân cận B, C; sau hai lớp, thông tin đạt đến các nút bậc 2 (lân cận của lân cận).
\end{tcolorbox}

\section{Mạng LSTM và Biến thể Stacked}

\subsection{Cơ chế Gate của LSTM}

Mạng LSTM \cite{hochreiter1997lstm} giải quyết vấn đề Vanishing Gradient của RNN thông thường thông qua cơ chế \textbf{cổng (gate)} điều tiết luồng thông tin. Tại mỗi bước thời gian $t$:

\begin{align}
  \mathbf{f}_t &= \sigma(\mathbf{W}_f[\mathbf{h}_{t-1}, \mathbf{x}_t] + \mathbf{b}_f) \quad \text{(Forget Gate)} \label{eq:forget} \\
  \mathbf{i}_t &= \sigma(\mathbf{W}_i[\mathbf{h}_{t-1}, \mathbf{x}_t] + \mathbf{b}_i) \quad \text{(Input Gate)} \label{eq:input} \\
  \tilde{\mathbf{c}}_t &= \tanh(\mathbf{W}_c[\mathbf{h}_{t-1}, \mathbf{x}_t] + \mathbf{b}_c) \quad \text{(Candidate Cell)} \label{eq:candidate} \\
  \mathbf{c}_t &= \mathbf{f}_t \odot \mathbf{c}_{t-1} + \mathbf{i}_t \odot \tilde{\mathbf{c}}_t \quad \text{(Cell State Update)} \label{eq:cell} \\
  \mathbf{o}_t &= \sigma(\mathbf{W}_o[\mathbf{h}_{t-1}, \mathbf{x}_t] + \mathbf{b}_o) \quad \text{(Output Gate)} \label{eq:output_gate} \\
  \mathbf{h}_t &= \mathbf{o}_t \odot \tanh(\mathbf{c}_t) \quad \text{(Hidden State)} \label{eq:hidden}
\end{align}

trong đó $\odot$ là tích Hadamard (element-wise multiplication), $\sigma$ là hàm sigmoid, và $[\cdot, \cdot]$ là phép ghép nối (concatenation).

\subsection{Stacked LSTM và Khả năng Biểu diễn Bậc cao}

Khi xếp chồng nhiều lớp LSTM (Stacked LSTM), mỗi lớp học một cấp độ trừu tượng khác nhau của chuỗi thời gian:

\begin{table}[H]
\centering
\caption{Ý nghĩa phân cấp trong Stacked LSTM áp dụng cho giao thông}
\label{tab:stacked_lstm}
\renewcommand{\arraystretch}{1.4}
\begin{tabular}{>{\bfseries\centering}p{1.8cm} p{4cm} p{6.5cm}}
\toprule
\textbf{Lớp} & \textbf{Đặc trưng học được} & \textbf{Ví dụ cụ thể} \\
\midrule
LSTM-1 & Biến động ngắn hạn (5--15 phút) & Đột biến lưu lượng do đèn đỏ \\
\rowcolor{rowgray}
LSTM-2 & Xu hướng trung hạn (15--45 phút) & Sóng ùn tắc từ một sự cố lan ra \\
LSTM-3 & Chu kỳ giờ cao điểm (1--2 giờ) & Sáng đông xe -- Trưa thưa -- Chiều kẹt \\
\bottomrule
\end{tabular}
\end{table}

\section{Mô hình Thay thế Nơ-ron (Neural Surrogate Model)}

\subsection{Hạn chế của Mô phỏng Vi mô}

Các công cụ mô phỏng vi mô (Microscopic Simulation) như SUMO (Simulation of Urban MObility) hay Vissim cho phép mô phỏng từng phương tiện riêng lẻ với độ chính xác vật lý cao. Tuy nhiên, độ phức tạp tính toán là $\mathcal{O}(N^2)$ với $N$ là số phương tiện:

\begin{table}[H]
\centering
\caption{So sánh Mô phỏng Vi mô và Surrogate Model}
\label{tab:sim_comparison}
\renewcommand{\arraystretch}{1.4}
\begin{tabular}{l p{4cm} p{4cm}}
\toprule
\textbf{Tiêu chí} & \textbf{Mô phỏng Vi mô (SUMO)} & \textbf{Neural Surrogate (ADE)} \\
\midrule
Độ phức tạp & $\mathcal{O}(N^2)$ --- phụ thuộc quy mô & $\mathcal{O}(1)$ --- cố định \\
\rowcolor{rowgray}
Thời gian chạy & 10 phút -- vài giờ & 100--300 ms \\
Độ chính xác & Rất cao (mô phỏng vật lý) & Xấp xỉ (học từ dữ liệu lịch sử) \\
\rowcolor{rowgray}
Xử lý Corner Cases & Tốt (mô hình vật lý đầy đủ) & Yêu cầu ADE + RealGen \\
Khả năng Real-time & Không & Có \\
\bottomrule
\end{tabular}
\end{table}

\subsection{Lý thuyết Deep Ensemble và ADE}

Deep Ensemble \cite{lakshminarayanan2017ade} là kỹ thuật đơn giản nhưng hiệu quả để ước lượng độ bất định (uncertainty):

\begin{equation}
  \hat{y}_{\text{ensemble}} = \frac{1}{M}\sum_{m=1}^{M} f_{\theta_m}(\mathbf{x}), \qquad
  \text{Var}[\hat{y}] = \frac{1}{M}\sum_{m=1}^{M}\left(f_{\theta_m}(\mathbf{x}) - \hat{y}_{\text{ensemble}}\right)^2
  \label{eq:ensemble}
\end{equation}

ADE (Adversarial Diverse Deep Ensemble) bổ sung hai cơ chế:
\begin{enumerate}[itemsep=3pt]
  \item \textbf{Diversity:} Các sub-model $f_{\theta_1}, \ldots, f_{\theta_M}$ có kiến trúc khác nhau (MLP, CNN-1D, Transformer nhẹ) để tăng đa dạng dự báo.
  \item \textbf{Adversarial Training:} Huấn luyện với các mẫu đối kháng $\mathbf{x}' = \mathbf{x} + \epsilon \cdot \nabla_\mathbf{x} \mathcal{L}$ (theo FGSM) để mô hình bền vững với các kịch bản cực trị (Corner Cases).
\end{enumerate}

\section{Truy xuất Tăng cường Phát sinh (RAG)}

\subsection{Kiến trúc RAG Cổ điển}

RAG \cite{lewis2020rag} là kỹ thuật kết hợp LLM với nguồn tri thức ngoại vi để giảm thiểu ảo giác (hallucination). Quy trình gồm hai giai đoạn:

\begin{enumerate}[itemsep=4pt]
  \item \textbf{Retrieval (Truy xuất):} Cho câu truy vấn $q$, tính vector nhúng $\mathbf{e}_q = \text{Enc}(q)$, sau đó tìm $k$ đoạn văn bản $\{d_1, \ldots, d_k\}$ có độ tương đồng cosine cao nhất trong Vector Database:
    \begin{equation}
      \text{sim}(\mathbf{e}_q, \mathbf{e}_d) = \frac{\mathbf{e}_q \cdot \mathbf{e}_d}{\|\mathbf{e}_q\|\|\mathbf{e}_d\|}
      \label{eq:cosine}
    \end{equation}
  \item \textbf{Generation (Sinh văn bản):} LLM sinh câu trả lời dựa trên ngữ cảnh được bổ sung: $p(y|q, d_1, \ldots, d_k)$.
\end{enumerate}

\subsection{Schema-Level RAG và Text-to-SQL}

RAG truyền thống chỉ xử lý văn bản phi cấu trúc. Trong STWI, dữ liệu số liệu mô phỏng có cấu trúc bảng (structured data) không thể truy vấn bằng semantic search. \textbf{Schema-Level RAG} giải quyết vấn đề này bằng cách chuyển đổi câu hỏi tự nhiên thành câu lệnh SQL thông qua mô hình Text-to-SQL (XiYanSQL \cite{guo2023xiyansql}).

\section{Hệ thống Đa Tác tử (Multi-Agent Systems)}

\subsection{Định nghĩa và Phân loại Tác tử}

Một \textbf{tác tử} (agent) là thực thể tính toán có khả năng: \textit{nhận thức môi trường} (perception), \textit{suy luận} (reasoning) và \textit{hành động} (action) \cite{wooldridge2009multiagent}. Trong STWI, mỗi tác tử AI là một LLM được trang bị bộ công cụ (tools):

\begin{table}[H]
\centering
\caption{Phân loại và đặc điểm các Tác tử trong STWI}
\label{tab:agents}
\renewcommand{\arraystretch}{1.4}
\begin{tabular}{p{3.5cm} p{3.5cm} p{5.5cm}}
\toprule
\textbf{Tác tử} & \textbf{Công cụ (Tools)} & \textbf{Nhiệm vụ} \\
\midrule
Simulation Agent & API Tầng 2 (ADE) & Chạy mô phỏng What-if \\
\rowcolor{rowgray}
Legal \& SOP Agent & Vector DB, XiYanSQL & Tra cứu pháp lý, truy vấn số liệu \\
Evaluation Agent & Tổng hợp kết quả & Đánh giá và đề xuất phương án \\
\rowcolor{rowgray}
Orchestrator & Điều phối 3 agent trên & Phân rã nhiệm vụ, hợp nhất kết quả \\
\bottomrule
\end{tabular}
\end{table}

\subsection{Kiến trúc ReAct và Chain-of-Thought}

Các tác tử trong STWI hoạt động theo kiến trúc \textbf{ReAct} (Reasoning + Acting): tại mỗi bước, tác tử luân phiên thực hiện \textit{Thought} (suy nghĩ), \textit{Action} (gọi tool), và \textit{Observation} (quan sát kết quả) cho đến khi đạt được mục tiêu.

\section{Suy luận Phản thực tế (Counterfactual Reasoning)}

\subsection{Lý thuyết Phản thực tế của Pearl}

Suy luận nhân quả (causal reasoning) theo khung của Pearl \cite{pearl2018why} phân biệt ba tầng:

\begin{enumerate}[leftmargin=1.5cm, itemsep=4pt]
  \item \textbf{Tương quan (Association):} ``Những ngày mưa, lưu lượng tại nút A thấp hơn.''
  \item \textbf{Can thiệp (Intervention):} ``Nếu tôi đóng đường X [do-calculus], lưu lượng ở Y sẽ thay đổi thế nào?''
  \item \textbf{Phản thực tế (Counterfactual):} ``Nếu thay vào đó tôi đóng đường Z (thay vì X), hậu quả sẽ tệ hơn hay tốt hơn?''
\end{enumerate}

\subsection{CF-VLA trong STWI}

CF-VLA (Counterfactual--Vision--Language--Action) trong STWI thực hiện tầng suy luận thứ 3: sau khi sinh một đề xuất sơ bộ, hệ thống \textit{tự đặt câu hỏi phản thực tế}: ``Nếu thực hiện hành động này, các nút giao lân cận có bị quá tải không?'' Câu trả lời được xác minh bằng mô phỏng ADE lần thứ hai, trước khi đề xuất được phê duyệt.

\begin{tcolorbox}[dangerbox]
\textbf{Tại sao CF-VLA là bắt buộc?} Nhiều hệ thống AI giao thông hiện tại đề xuất phân luồng mà không kiểm tra tác động bậc 2. Ví dụ: phân luồng toàn bộ xe từ đường A sang đường B có thể giải quyết ùn tắc tại A nhưng gây ùn tắc nghiêm trọng hơn tại ngã tư C (nút giao của đường B với đường khác). CF-VLA phát hiện và ngăn chặn kịch bản này.
\end{tcolorbox}

\section{Độ Đo Đánh giá Hệ thống}

Các độ đo chính được sử dụng xuyên suốt trong báo cáo:

\begin{align}
  \text{RMSE} &= \sqrt{\frac{1}{n}\sum_{i=1}^n (\hat{y}_i - y_i)^2} \label{eq:rmse} \\
  \text{MAE}  &= \frac{1}{n}\sum_{i=1}^n |\hat{y}_i - y_i| \label{eq:mae} \\
  \text{F1}   &= 2 \cdot \frac{\text{Precision} \times \text{Recall}}{\text{Precision} + \text{Recall}} \label{eq:f1} \\
  \text{Recall@}k &= \frac{|\text{Rel}_k \cap \text{Ret}_k|}{|\text{Rel}|} \label{eq:recall_k}
\end{align}

Trong đó $\hat{y}_i$ là giá trị dự báo, $y_i$ là giá trị thực tế, $\text{Rel}_k$ là tập tài liệu liên quan trong top-$k$ kết quả truy xuất, $\text{Ret}_k$ là tập tài liệu được truy xuất.
